\oldpage[IV-2]{176}
V) \emph{Fin de la demostración.} --- Por II), se puede suponer que
$B$ es una $A$-álgebra de tipo finito que es un dominio entero y
contiene a $A$. Sean $K$ el cuerpo de fracciones de $A$ y $L$ el de
$B$, que es una extensión de tipo finito de $K$. Sea $M$ la
clausura algebraica de $K$ en $L$, una extensión algebraica \emph{finita} de $K$; sea $C_0$ la clausura integral de $A$ en $M$, que es también
la clausura integral de $A$ en $L$. Se tiene, pues, $C=B\cap C_0$.
Si se pone $C'_0=C\otimes_A A'$, se sigue por planitud que
$C'=B'\cap C'_0$ ($\mathbf{0}_{\mathrm I}$, 6.1.3). Por IV), $C'_0$ es la
clausura integral de $A'$ en $M'=M\otimes_A A'$. Además, $M'$ es un
anillo noetheriano y el morfismo
$\operatorname{Spec}(M')\to\operatorname{Spec}(M)$ es normal, como se
vio en IV). Aplicando III) a $M$ y $M'$ en lugar de $A$ y $A'$, y
a $L$ en lugar de $B$, se deduce que $M'$ está integralmente cerrado
en $L'=L\otimes_A A'$. Por tanto, $C'_0$ es la clausura integral de
$A'$ en $L'$, y $C'=C'_0\cap B'$ es la clausura integral de $A'$ en
$B'$. Q.E.D.
\end{proof}

